\documentclass{article}
\usepackage{/home/user1/code/tex/sty}
\usepackage{amsfonts}
\usepackage{amsmath}
\usepackage{geometry}
\usepackage{graphicx}
\usepackage{hyperref}
\usepackage{enumitem}
\setlength{\parindent}{0pt}
\geometry{top=1in,bottom=1in,left=1in,right=1in}
\begin{document}
\begin{center}
\textbf{\Large{Woche 3 - Hausafgaben}}\\~\\
\begin{tabular}{l l}
Student&Agustin Romero\\
Email&ar1@stanford.edu\\
Instructor&Patric Di Dio De Marco\\
Course&GERLANG 3\\
Institution&Stanford University\\
Date&\today\\
\end{tabular}
\end{center}
\section*{Hausafgaben Woche 3 Tag 1}
\subsection*{181-5}
Schreiben Sie die Komparative.
\begin{itemize}
\item Die Komparativ von ,,nah'' ist näher.
\item Die Komparativ von ,,warm'' ist wärmer.
\item Die Komparativ von ,,schnell'' ist scheller.
\item Die Komparativ von ,,ruhig'' ist ruhiger.
\item Die Komparativ von ,,billig'' ist billiger.
\item Die Komparativ von ,,kurz'' ist kürzer.
\end{itemize}
\begin{enumerate}[label=(\alph*)]
\item Der Flug nach New York kostet 800 €. 
\begin{itemize}
\item Das ist zu teuer. Wir müssen einen \underline{billigeren} Flug finden.
\end{itemize}
\item Der flug auf die Philippinen dauert 15 Stunden. Ist das zu lang?
\begin{itemize}
\item Ja, such wir doch ein \underline{kurzer} Urlaubsziel.
\end{itemize}
\item Mit unseren alten Fahrrädern dauert die Tour mindesetens drei Tage.
\ite{\item Dann leihen wir uns doch \underline{schneller} Räder.}
\item Willst du über Hamburg oder über Leipzig nach Berlin fahren?
\ite{\item Ich bin für die \underline{kurzer} Strecke.}
\item Ich möchte im Sommer nach Grönland reisen.
\ite{\item Ich würde lieber in ein \underline{warmer} Land fahren.}
\item Fahren wir doch eine Woche nach Berlin.
\ite{\item Ein \underline{näher} Ort wäre mir lieber.}
\end{enumerate}
\subsection*{181-6}
\begin{enumerate}[label={\alph*}]
\item Onkel Michael meint, früher gab es \underline{gringere Verkehrprobleme}.
\item Lilli meint, dass es heute \underline{preiswertere U-Bahnen gibt}.
\item Michael findet, dass es \underline{früher vernünftigere Freizeitbeschäftigungen}.
\item Lilli gibt ihrem Onkel den Rat, einen Computer zu kaufen, weil es keine \underline{einfachere Kommunikationsmöglichkeit gibt}.
\item Lilli findet, im Internet gibt es oft \underline{niedrigere Preise} als im Supermarkt.
\item Michael fragt, warum es \underline{zuverlässigere Verkehrsmittel} als grüher gibt, obwohl jeder einen Computer hat.
\end{enumerate}
\section*{Hausafgaben Woche 3 Tag 2}
\subsection*{182-2}
\itealph{\item Wir könnten \und{irgendwo} auf einem Parkplatz schlafen.
\ite{\item Ich will aber in einem guten Hotel schlafen.}
\item Ich möchte gern im Urlaub zu Hause bleiben. Müssen wir immer \und{irgendwohin} fahren?
\ite{\item Ja, ich will unbedingt nach Frankreich.}
\item Wil sollten \und{irgendetwas} als Geschenk mitbringen.
\ite{\item Stimmt. Vielleicht eine Flasche Rotwein?}
\item Wir können uns ja \und{irgendwann} nächste Woche treffen.
\ite{\item Gute Idee! Wir wär's mit Donnerstag?}
\item Ich glaube, wir sollten \und{irgendwen} nach dem Weg fragen.
\ite{\item Fragen wir doch den Mann dort drüben.}
\item \und{Irgendwie} müssen wir jetzt nach Frankfurt kommen.
\ite{\item Fahren wir doch per Autostopp.}}
\subsection*{183-5}
\section*{Hausafgaben Woche 3 Tag 3}
\subsection*{183-6}
\subsection*{183-8}
\subsection*{183-9}
\subsection*{184-4}
\section*{Hausafgaben Woche 3 Tag 4}
\end{document}
